\section*{Overview}
\label{sec:overview}
\addoverviewtotoc

This blueprint follows the order of the paper.  It first develops Novikov
series and Novikov isocrystals, then introduces Novikov descent data and the
comparison between descent data and isocrystals.  The remaining sections prove
essential surjectivity by passing successively through algebraically closed
fields, products of such fields, reduced rings, arbitrary coefficient rings,
and finally arbitrary exponent monoids.

\paragraph{Numbering convention.}
Statements are numbered locally within each blueprint section.  A
parenthetical label such as ``Proposition~4.8'' records the corresponding
number in the paper; headings without such a label name supporting results
introduced for the formalization.

\paragraph{Main theorem.}
The main result is Theorem~1.1 of the paper, restated and proved there as
Theorem~4.9.  Let $A$ be a commutative ring and let
$\Gamma\subseteq\mathbb{R}$ be an additive submonoid.  The goal is the equivalence
\[
  \operatorname{Const}_\Gamma\colon
  \mathsf{Vect}(A)\xrightarrow{\ \sim\ }\mathsf{NovDD}(A,\Gamma),
\]
where the source is the category of finite projective $A$-modules and the
target is the category of Novikov descent data.  This is proved in
Theorem~\ref{thm:novikov_descent_equivalence}.

\paragraph{Proof architecture.}
The argument has three stages.
\begin{enumerate}
  \item The sections on Novikov series, isocrystals, and descent establish the
    algebraic foundations, full faithfulness of constant objects, and the
    fully faithful comparison from real-exponent descent data to isocrystals.
  \item For real exponents, essential surjectivity is proved first over an
    algebraically closed field, then over products of algebraically closed
    fields and reduced rings.  Square-zero deformation and a finitely
    generated nilpotent quotient reduce the general coefficient ring to the
    reduced case.
  \item For a general $\Gamma$, exponents are extended to $\mathbb{R}$.  The
    real-exponent theorem supplies a trivialization, support arguments restrict
    it back to $\Gamma$, and finite-projective duality upgrades the resulting
    injection to an isomorphism.
\end{enumerate}

\paragraph{Notation.}
Throughout the blueprint:
\begin{description}
  \item[$A((t_1^\Gamma,\dotsc,t_n^\Gamma))$] denotes the $n$-variable Novikov
    ring with coefficients in $A$ and exponents in $\Gamma$; in one variable
    we write $A((t^\Gamma))$.
  \item[$\mathsf{Vect}(A)$] denotes the category of finite projective
    $A$-modules.
  \item[$\mathsf{NovIsoc}(A)$] denotes the category of real-exponent Novikov
    isocrystals for the fixed dilation factor $\lambda>1$.
  \item[$\mathsf{NovDD}(A,\Gamma)$] denotes the category of Novikov descent
    data with exponent monoid $\Gamma$.  When $\Gamma=\mathbb{R}$, the exponent
    is sometimes omitted.
  \item[$\operatorname{Const}_\Gamma(P)$] denotes the constant descent datum
    associated to $P$; the subscript is omitted in the real-exponent sections.
  \item[$M_B$] denotes coefficient base change along a specified map
    $A\to B$, while $M_{\mathbb{R}}$ denotes extension of exponents from
    $\Gamma$ to $\mathbb{R}$.
  \item[$\pi_i^\ast$] denotes a face map and, by the same notation, the induced
    scalar extension of modules along that map.
\end{description}
